\documentclass[11pt,a4paper]{article}
\usepackage[margin=1in]{geometry}
\usepackage{amsmath,amssymb}
\usepackage{graphicx}
\usepackage{float}
\usepackage{enumitem}
\usepackage{microtype}
\usepackage{caption}
\usepackage{tabularx}
\usepackage{hyperref}

\captionsetup{font=small,labelfont=bf,skip=6pt}

\setlist[itemize]{leftmargin=1.2cm}
\setlist[enumerate]{leftmargin=1.2cm}

\newcommand{\problembox}[3]{%
  \par\addvspace{0.42em}%
  \noindent
  \setlength{\fboxsep}{7.5pt}%
  \setlength{\fboxrule}{0.95pt}%
  \fbox{%
    \begin{minipage}{\dimexpr\linewidth-2\fboxsep-2\fboxrule\relax}
    \begin{tabularx}{\linewidth}{@{}>{\raggedright\arraybackslash\bfseries}p{2.1em} @{\hspace{0.9em}} >{\raggedright\arraybackslash}X @{\hspace{0.9em}} >{\raggedleft\arraybackslash\bfseries}p{2.4em}@{}}
    #1 & #2 & #3
    \end{tabularx}
    \end{minipage}%
  }%
  \par\addvspace{0.32em}%
}


\title{\textbf{Electron tube}}
\author{\normalsize Y25 (2025) --- English translation}
\date{}
\begin{document}
\maketitle

This problem is devoted to a device called an electron tube. Electron tubes differ in shape and design, but the principle of operation is the same: in an evacuated flask, two electrodes are located: a cathode and an anode. The cathode is heated by an electric current flowing through it (or a nearby heater), leading to thermionic emission of electrons from its surface. A potential difference between the cathode and anode drives the movement of electrons in the vacuum. Other elements, such as grids, may also be present.

In all answers, unless otherwise specified, you may use:
\begin{itemize}
    \item magnetic constant $\mu_0 = 4\pi \cdot 10^{-7}$ H/A$^2$;
    \item electrical constant $\varepsilon_0 = 8.85 \cdot 10^{-12}$ F/m;
    \item Boltzmann constant $k = 1.38 \cdot 10^{-23}$ J/K;
    \item electron mass $m = 9.1 \cdot 10^{-31}$ kg;
    \item electron charge magnitude $e = 1.6 \cdot 10^{-19}$ C.
\end{itemize}

\textit{Note: The problem will require solving equations of the form $\psi''(z) = f(\psi(z))$. Multiplying both sides by $\psi'(z)$ and integrating yields $\frac{(\psi'(z))^2}{2} = \int f(\psi) d\psi + C$, where $C$ is a constant determined by initial conditions.}

\section*{Part A. Motion of an electron in an electromagnetic field (3.8 points)}

Consider a vacuum diode where the cathode and anode are parallel plates of area $S$. The distance between them is $h \ll \sqrt{S}$. The cathode is heated by a straight long wire of radius $r_0$ pressed against it, carrying a current $I$. Let the $x$-axis be parallel to the plates, the $y$-axis perpendicular to them (upward), and the $z$-axis directed towards us. The origin $C$ is above the wire.

\problembox{A1}{Let an electron start with zero initial speed from point $C$. Show that $x(t) = 0$.}{0.30}

\problembox{A2}{Write Newton's second law for the $y$ and $z$ projections and find expressions for $\ddot{y}$ and $\ddot{z}$. These may contain $\dot{y}, \dot{z}, y, z, U, I$ and system geometry.}{0.50}

\problembox{A3}{Obtain a differential equation for $y$ of the form: $\dot{y}^2 = \alpha - \beta \frac{\ln(1+y/r_0)}{1+y/r_0}$. Find the constants $\alpha$ and $\beta$.}{1.00}

\problembox{A4}{Find the maximum heater current $I_{cr}$ at which the electron reaches the anode.}{0.80}

Now consider a cylindrical configuration: the cathode is a wire of radius $r_0$ and length $L$ carrying current $I$; the anode is a coaxial cylinder of radius $R_0$ and length $L$ ($L \gg R_0$). The potential difference is $U = \varphi_A - \varphi_K > 0$.

\problembox{A5}{Find the maximum heater current $I_{cr}$ for which an electron starting from point $C$ with zero speed reaches the anode in this cylindrical configuration.}{1.20}

\section*{Part B. I-V characteristics of a vacuum diode (4.1 points)}

Assume electrons move in a steady state, creating a spatial distribution of charge $\rho$ and potential $\varphi$. We consider only the electric field ($I \ll I_{cr}$) and zero initial electron velocity ($kT \ll eU$). Let $\varphi(0) = 0$ at the cathode ($y=0$).

\problembox{B1}{Derive a differential equation relating $\varphi(y)$ and $\rho(y)$.}{0.30}

\problembox{B2}{For $U > 0$, find the electron speed $v(y)$ at distance $y$ from the cathode. Express it through $\varphi(y)$.}{0.20}

\problembox{B3}{Express the total current $I$ (positive from anode to cathode) through $\rho(y), v(y)$, and system geometry.}{0.30}

Assume the electric field at the cathode is zero ($\varphi'(0) = 0$).

\problembox{B4}{Obtain a differential equation for $\varphi(y)$ containing $I$ and system geometry.}{0.40}

\problembox{B5}{Find the dependence $\varphi(y)$ in terms of $I$ and system geometry.}{0.80}

\problembox{B6}{The I-V characteristic is $I(U) = GU^{\gamma}$. Determine the perveance $G$ and exponent $\gamma$.}{0.80}

\problembox{B7}{Find the charge density distribution $\rho(y)$.}{0.50}

\problembox{B8}{Determine the I-V characteristic $I(U)$ for $U < 0$.}{0.20}

\problembox{B9}{For a cylindrical configuration with radial coordinate $r$, obtain the differential equation for $\varphi(r)$ involving total current $I$ and geometry.}{0.60}

\section*{Part C. Triode and Van der Pol Oscillator (2.1 points)}

A conducting grid with potential $\varphi_C$ is placed between the plates. $U_C = \varphi_C - \varphi_K$ is grid voltage; $U_A = \varphi_A - \varphi_K$ is anode voltage. Assume no current flows to the grid. The anode current $I_A$ depends on $U_C$ and $U_A$. For $U_A = \text{const}$ near $U_C = 0$:
$$I_A(U_C) \simeq I_0 + \lambda U_C - \mu U_C^3 \quad (1)$$
where $\lambda, \mu$ are positive constants.

\begin{figure}[H]
    \centering
    \includegraphics[width=0.45\textwidth]{T6_EN_assets/fig6.jpg}
    \caption{Van der Pol oscillator circuit.}
\end{figure}

Assume the EMF induced by coil $L$ in $L_A$ is negligible. The coil $L_A$ induces an EMF in $L$: $\mathcal{E} = M \frac{dI_A}{dt}$, where $M$ is the mutual inductance.

\problembox{C1}{Using Kirchhoff's law for the RLC circuit, derive an equation: $\ddot{U}_C + \zeta(1+\eta U_C^2)\dot{U}_C + \chi U_C = 0$. Express $\zeta, \eta, \chi$ through $R, L, C, M, \lambda, \mu$.}{0.90}

\problembox{C2}{Find the condition on $\zeta, \eta, \chi$ for automatic oscillation buildup from $U_C \approx 0$. Express this condition in terms of $R, L, C, M, \lambda, \mu$.}{1.20}

\vspace{1em}
\noindent\textit{Source:} https://pho.rs/p/4223

\end{document}